\subsection{Cone-compactified sphere carriers for the native cusp quotient}
\label{cusp-cone-section}

This subsection answers a topological question which is strictly different
from the finite-cover construction.  No period is enlarged by passing to a
cover.  Instead, the native expanding quotient is retained on the cusp, the
parts on which that quotient is unavailable are collapsed, and the resulting
based map is passed through a smash product.  The outcome is an exact global
map with source \(S^{12}\), and there is an analogous \(S^{24}\) map from the
full-monodromy-invariant circle quotient.  The targets are conically
stratified spaces.  The Euclidean limits below use the intrinsic metrics
inherited from the flat quotient tori on the regular slices; they are not
ambient metric limits in one fixed metric
on the compact targets.

Fix a closed coordinate disc
\begin{equation}
 D_0=\{t\in\mathbb C:|t|\leq\epsilon\}\subset B,\qquad
 p_0=\{t=0\},\qquad N_0=f^{-1}(D_0),\qquad W=f^{-1}(0),
 \label{cusp-cone-disc}
\end{equation}
small enough that it contains no singular value other than \(p_0\).  Put
\begin{equation}
 K_\Lambda=\mathbb Z\langle\widehat w,\widehat\delta\rangle,\qquad
 \overline\Lambda=\Lambda/K_\Lambda,\qquad
 \overline T^2=
 \frac{\overline\Lambda\otimes\mathbb R}{\overline\Lambda}.
 \label{cusp-cone-quotient-lattice}
\end{equation}
The ordered basis of \(\overline\Lambda\) is the image of
\((\widehat\gamma,\widehat u)\).

\begin{lemma}[The native quotient over the whole punctured cusp]
\label{cusp-cone-native-descent}
There is a well-defined real-analytic fibrewise group epimorphism
\begin{equation}
 \overline\pi:N_0\setminus W
 \longrightarrow D_0^*\times\overline T^2
 \label{cusp-cone-native-family-map}
\end{equation}
whose kernel is the connected two-torus family defined by \(K_\Lambda\).
The target local system is trivial.  On the positive radial fibre \(F_T\),
\(t=e^{-2\pi T}\), the isomorphism
\begin{equation}
 \iota_T:\overline T^2\longrightarrow Q_T,\qquad
 [(a,b)]\longmapsto[Y_T(a,b)^{\mathsf T}]
 \label{cusp-cone-radial-isometry}
\end{equation}
satisfies
\begin{equation}
 \pi_T=\iota_T\circ\overline\pi|_{F_T}.
 \label{cusp-cone-native-restriction}
\end{equation}
If \(\overline T^2\) is given the pullback by \(\iota_T\) of the
Euclidean quotient metric on \(Q_T=\mathbb R^2/Y_T\mathbb Z^2\), then
\(\iota_T\) is an isometry and
\begin{equation}
 \operatorname{sys}(\overline T^2,g_T)
 =\operatorname{sys}(Q_T)\geq T-C.
 \label{cusp-cone-native-systole}
\end{equation}
\end{lemma}

\begin{proof}
On the universal cover of \(D_0^*\), use the real-analytic flat marking
\[
 (z,x+\Lambda)\longmapsto[z,\Pi(z)x].
\]
The deck generator acts by
\[
 (z,x+\Lambda)\longmapsto(g_0z,M_0x+\Lambda).
\]
The displayed matrix \(M_0\) gives
\begin{equation}
 (M_0-I)(a,b,c,d)^{\mathsf T}=(0,0,b,-a)^{\mathsf T},\qquad
 \operatorname{im}(M_0-I)=K_\Lambda.
 \label{cusp-cone-monodromy-image}
\end{equation}
Consequently \(M_0\) preserves \(K_\Lambda\) and induces the identity on
\(\overline\Lambda\).  The formula
\[
 [z,\Pi(z)x]\longmapsto
 \bigl([z],[x\bmod K_\Lambda]\bigr)
\]
is invariant under the deck action and proves
\cref{cusp-cone-native-family-map}.  The lattice \(K_\Lambda\) is primitive,
so the kernel is connected.  On the radial fibre, projection to the
\((e_2,e_4)\)-coordinates sends
\(a\widehat\gamma+b\widehat u\) to \(Y_T(a,b)^{\mathsf T}\) and kills
\(K_\Lambda\).  This proves
\cref{cusp-cone-radial-isometry,cusp-cone-native-restriction}.  The metric
and systolic assertion is the quotient calculation in
\cref{cusp-spec-native-quotient}.
\end{proof}

Define
\begin{equation}
 \mathcal A=
 \frac{[0,\epsilon]\times S^1\times\overline T^2}
 {\{0\}\times S^1\times\overline T^2\sim v_0,\quad
  \{\epsilon\}\times S^1\times\overline T^2\sim v_1}.
 \label{cusp-cone-A-definition}
\end{equation}
Thus
\begin{equation}
 \mathcal A\cong
 \operatorname{Susp}(S^1\times\overline T^2)
 \cong\operatorname{Susp}(T^3),
 \label{cusp-cone-A-suspension}
\end{equation}
where \(\operatorname{Susp}\) denotes the unreduced suspension.  The extra
circle is the angular coordinate of the punctured base.  If one deliberately
forgets that angle, the evident map
\(\mathcal A\to\operatorname{Susp}(\overline T^2)\) gives a smaller target
with the same restriction to every fixed radial fibre.  We retain
\(\mathcal A\) because it records the entire punctured-cusp quotient.

\begin{theorem}[Exact localized \(S^{12}\) route]
\label{cusp-cone-S12-route}
Choose \(x_0\in W\), and base \(\mathcal A\) at \(v_0\).  There is a
continuous surjective based quotient map
\begin{equation}
 q_2:(X,x_0)\longrightarrow(\mathcal A,v_0)
 \label{cusp-cone-q2}
\end{equation}
with the following exact description.  For \(x\in N_0\setminus W\), write
\[
 t(f(x))=re^{i\phi},\qquad 0<r<\epsilon,
\]
and let \(\overline\pi_f(x)\in\overline T^2\) be the fibre component of
\(\overline\pi(x)\).  Then
\begin{equation}
 q_2(x)=[r,e^{i\phi},\overline\pi_f(x)],\qquad
 q_2(W)=v_0,\qquad
 q_2\bigl(X\setminus\operatorname{int}N_0\bigr)=v_1.
 \label{cusp-cone-q2-formula}
\end{equation}
On \(f^{-1}(|t|=\epsilon)\), the value is \(v_1\).

The smash square is a global based map
\begin{equation}
 q_2^{\wedge2}:X\wedge X\longrightarrow
 \mathcal A\wedge\mathcal A.
 \label{cusp-cone-S12-map}
\end{equation}
Using \(X\cong_{\mathrm{diff}}S^6\), a based choice of the diffeomorphism
identifies its source with
\begin{equation}
 X\wedge X\cong S^6\wedge S^6\cong S^{12}.
 \label{cusp-cone-S12-source}
\end{equation}
For \(0<e^{-2\pi T}<\epsilon\), the canonical product-to-smash quotient
maps restrict to embeddings on \(F_T^2\) and on the corresponding interior
target-product cross-section.  Under those embeddings and \(\iota_T^2\),
the induced smash map restricts exactly to
\begin{equation}
 \pi_T\times\pi_T:F_T\times F_T\longrightarrow Q_T\times Q_T.
 \label{cusp-cone-S12-fibre-map}
\end{equation}
With the Euclidean \(\ell^2\)-product flat quotient metric,
\begin{equation}
 \operatorname{sys}(Q_T\times Q_T)\geq T-C,\qquad
 (Q_T\times Q_T,(0,0))\longrightarrow(\mathbb R^4,0)
 \label{cusp-cone-S12-R4-limit}
\end{equation}
in pointed local-isometric, pointed smooth, and pointed
Gromov--Hausdorff topology.  Once \(T-C>4R\), the radius-\(R\) ball is
exactly Euclidean.

The integral homology of the targets is
\begin{equation}
 \widetilde H_k(\mathcal A;\mathbb Z)\cong
 \begin{cases}
  \mathbb Z^3,&k=2,3,\\
  \mathbb Z,&k=4,\\
  0,&\text{otherwise},
 \end{cases}
 \label{cusp-cone-A-homology}
\end{equation}
and
\begin{equation}
 \widetilde H_k(\mathcal A\wedge\mathcal A;\mathbb Z)\cong
 \begin{cases}
  \mathbb Z^9,&k=4,\\
  \mathbb Z^{18},&k=5,\\
  \mathbb Z^{15},&k=6,\\
  \mathbb Z^6,&k=7,\\
  \mathbb Z,&k=8,\\
  0,&\text{otherwise}.
 \end{cases}
 \label{cusp-cone-A-smash-homology}
\end{equation}
In particular, the target is not \(S^4\), \(S^{12}\), or a compact
four-manifold.
\end{theorem}

\begin{proof}
The equivalence relation in \cref{cusp-cone-A-definition} is closed: it is
the union of the diagonal with the squares of the two closed end
cross-sections.  A compact Hausdorff space modulo a closed equivalence
relation is compact Hausdorff.  Rescaling \(r\) to the unit interval gives
exactly the unreduced suspension in \cref{cusp-cone-A-suspension}.

Well-definedness of the interior formula follows from
\cref{cusp-cone-native-descent}.  A full turn in \(\phi\) changes the
marking by \(M_0\), but \(M_0\) is the identity modulo \(K_\Lambda\).

We check both collapsed strata.  If \(U\) is an open neighborhood of
\(v_0\), its inverse image in the cylinder is open and contains the compact
set \(\{0\}\times S^1\times\overline T^2\).  Hence, for some \(\delta>0\),
it contains
\([0,\delta)\times S^1\times\overline T^2\).  Therefore
\(q_2^{-1}(U)\) contains \(f^{-1}(|t|<\delta)\), a neighborhood of \(W\).
The same compactness argument at the other end shows that every neighborhood
of \(v_1\) contains the image of
\((\epsilon-\delta,\epsilon]\times S^1\times\overline T^2\).  Its preimage
together with \(X\setminus\operatorname{int}N_0\) is a neighborhood of the
outer boundary.  Continuity on the interior is
\cref{cusp-cone-native-family-map}; the closed pasting lemma gives
continuity on \(X\).

Every interior point has a preimage because \(\overline\pi\) is fibrewise
surjective.  The nonempty sets \(W\) and
\(X\setminus\operatorname{int}N_0\) map onto \(v_0\) and \(v_1\),
respectively.  Thus \(q_2\) is surjective.  A continuous surjection from
compact \(X\) to Hausdorff \(\mathcal A\) is a quotient map.

The map is based, so functoriality gives \cref{cusp-cone-S12-map}.  The
standard CW quotient gives \(S^m\wedge S^n\cong S^{m+n}\)
\cite[Chapter~0, Smash Product]{Hatcher2002}, proving
\cref{cusp-cone-S12-source}.

The base point lies in \(W\), so \(F_T^2\) is disjoint from the source fat
wedge.  Its image lies in an interior cross-section of
\(\mathcal A^2\), disjoint from the target fat wedge.  The source and target
fat wedges are closed in the compact Hausdorff products \(X^2\) and
\(\mathcal A^2\), so both smash quotients are Hausdorff.  Their restrictions
to the relevant compact subspaces are continuous and one-to-one, and are
therefore embeddings.
Equations \cref{cusp-cone-native-restriction,cusp-cone-q2-formula} give
\cref{cusp-cone-S12-fibre-map}.

The period lattice of the product target is
\(\operatorname{diag}(Y_T,Y_T)\mathbb Z^4\).  Its systole is the systole of
\(Y_T\mathbb Z^2\), proving the first assertion in
\cref{cusp-cone-S12-R4-limit}.  If \(x,y\in B_R(0)\subset\mathbb R^4\)
and \(0\ne\lambda\) is a product period, then
\[
 \|x-y+\lambda\|\geq\|\lambda\|-\|x-y\|
 \geq(T-C)-2R>2R\geq\|x-y\|
\]
when \(T-C>4R\).  The quotient projection is an exact isometry on the
Euclidean radius-\(R\) ball.  Conversely, any point of quotient distance
less than \(R\) has a shortest lattice representative of norm less than
\(R\), so this embedded ball is the entire quotient radius-\(R\) ball.
This proves all three pointed convergences.

Reduced suspension shifts reduced homology by one.  Since
\[
 H_1(T^3)=\mathbb Z^3,\qquad
 H_2(T^3)=\mathbb Z^3,\qquad H_3(T^3)=\mathbb Z,
\]
\cref{cusp-cone-A-homology} follows.  All groups are free, so the Kunneth
theorem for smash products has no Tor terms.  The ranks in degrees
\(4,5,6,7,8\) are
\[
 3\cdot3,\quad 2(3\cdot3),\quad
 2(3\cdot1)+3\cdot3,\quad 2(3\cdot1),\quad1,
\]
which proves \cref{cusp-cone-A-smash-homology}.
\end{proof}

\paragraph{The localized one-circle construction.}
The full-monodromy-invariant character also gives a smaller cusp-supported
target.  Define
\begin{equation}
 \mathcal B=
 \frac{[0,\epsilon]\times S^1\times S^1}
 {\{0\}\times S^1\times S^1\sim b_0,\quad
  \{\epsilon\}\times S^1\times S^1\sim b_1}
 \cong\operatorname{Susp}(T^2),
 \label{cusp-cone-B-definition}
\end{equation}
where the two circle coordinates in an interior cross-section are the base
angle and the value of \(\theta_\gamma\).

\begin{proposition}[Localized \(S^{24}\) route with target
\(\operatorname{Susp}(T^2)^{\wedge4}\)]
\label{cusp-cone-S24-local-route}
There is a continuous surjective based quotient map
\begin{equation}
 q_{1,\mathrm{cusp}}:(X,x_0)\longrightarrow(\mathcal B,b_0)
 \label{cusp-cone-q1-local}
\end{equation}
given by
\begin{equation}
 q_{1,\mathrm{cusp}}(x)
 =[r,e^{i\phi},\theta_\gamma(x)]
 \quad(0<|t|=r<\epsilon),\qquad
 q_{1,\mathrm{cusp}}(W)=b_0,\qquad
 q_{1,\mathrm{cusp}}
 \bigl(X\setminus\operatorname{int}N_0\bigr)=b_1.
 \label{cusp-cone-q1-local-formula}
\end{equation}
Its fourth smash power has source \(X^{\wedge4}\cong S^{24}\), and on
four synchronized radial fibres it is exactly
\begin{equation}
 \chi_{\gamma,T}^{\times4}:F_T^4\longrightarrow C_T^4.
 \label{cusp-cone-S24-local-fibre-map}
\end{equation}
Consequently its intrinsic target slices satisfy the Euclidean limit
\cref{cusp-cone-S24-R4-limit}.  The target homology is
\begin{equation}
 \widetilde H_k(\mathcal B;\mathbb Z)\cong
 \begin{cases}
  \mathbb Z^2,&k=2,\\
  \mathbb Z,&k=3,\\
  0,&\text{otherwise},
 \end{cases}
 \label{cusp-cone-B-homology}
\end{equation}
and
\begin{equation}
 \widetilde H_k(\mathcal B^{\wedge4};\mathbb Z)\cong
 \begin{cases}
  \mathbb Z^{16},&k=8,\\
  \mathbb Z^{32},&k=9,\\
  \mathbb Z^{24},&k=10,\\
  \mathbb Z^8,&k=11,\\
  \mathbb Z,&k=12,\\
  0,&\text{otherwise}.
 \end{cases}
 \label{cusp-cone-B-smash-homology}
\end{equation}
\end{proposition}

\begin{proof}
The monodromy invariance
\(\varepsilon_\gamma M_0=\varepsilon_\gamma\) makes the interior formula
single-valued after a turn in the base angle.  The two endpoint-continuity
arguments are identical to those in the proof of
\cref{cusp-cone-S12-route}, with compact cross-section \(T^2\) in place of
\(T^3\).  Fibrewise surjectivity of \(\theta_\gamma\), together with the
nonempty inner fibre and exterior, proves surjectivity.  Compact-to-Hausdorff
makes the map a quotient map.

Since \(x_0\in W\), the fourth smash power exists.  The radial regular
fibres and their circle targets avoid the two closed fat wedges.  The
restrictions of the Hausdorff smash quotients to these compact subspaces are
continuous and one-to-one, hence embeddings, and the fibre map is exactly
\cref{cusp-cone-S24-local-fibre-map}.  The intrinsic metric statement is
therefore the four-circle calculation
\cref{cusp-cone-S24-R4-limit}.  Finally, suspension shifts the free homology
of \(T^2\), giving \cref{cusp-cone-B-homology}.  In the fourth smash power,
choosing \(j\) of the four factors in degree \(3\) and the other
\(4-j\) factors in degree \(2\) gives degree \(8+j\) and rank
\(\binom4j2^{4-j}\).  This is \cref{cusp-cone-B-smash-homology}.
\end{proof}

\paragraph{The full-base circle target before coning.}
Choose pairwise disjoint closed coordinate discs \(D_0,D_1,D_2\) about the
three singular values and put
\begin{equation}
 P=B\setminus\bigcup_{j=0}^2\operatorname{int}D_j.
 \label{cusp-cone-pair-of-pants}
\end{equation}
For \(j=1,2\), let \(C_j\) be the solid-torus Seifert target in
\cref{cusp-spec-finite-Seifert-target}, and let \(h_j^{S^1}\) be the
target overlap map in \cref{cusp-spec-finite-overlap-identity}.  Glue
\begin{equation}
 \mathcal M=(P\times S^1)
 \underset{h_1^{S^1}}{\cup}C_1
 \underset{h_2^{S^1}}{\cup}C_2.
 \label{cusp-cone-M-definition}
\end{equation}
The gluing is along boundary tori over \(\partial D_1,\partial D_2\).
The remaining boundary is
\begin{equation}
 \partial\mathcal M=\partial D_0\times S^1\cong T^2.
 \label{cusp-cone-M-boundary}
\end{equation}
Define
\begin{equation}
 \mathcal C=\mathcal M\cup_{\partial\mathcal M}
 \operatorname{Cone}(\partial\mathcal M),
 \label{cusp-cone-C-definition}
\end{equation}
and write \(c_0\) for the cone vertex.
This changes the codomain at the cusp: the whole boundary two-torus,
including the target circle, is collapsed to \(c_0\).  It therefore does
not contradict the proved nonextension of the literal map
\(\chi_\gamma\) to the fixed product target \(B\times S^1\).

\begin{theorem}[Exact cone-compactified \(S^{24}\) route]
\label{cusp-cone-S24-route}
There is a continuous surjective based quotient map
\begin{equation}
 q_1:(X,x_0)\longrightarrow(\mathcal C,c_0),\qquad x_0\in W,
 \label{cusp-cone-q1}
\end{equation}
defined as follows:
\begin{enumerate}
\item over \(P\), it is
\((f,\theta_\gamma):f^{-1}(P)\to P\times S^1\), equivalently the
restriction of \(\chi_\gamma\) to \(f^{-1}(P)\);
\item over \(D_j\), \(j=1,2\), it is the finite-filling map
\(N_j\to C_j\) from \cref{cusp-spec-finite-Seifert-target};
\item on \(N_0\setminus W\), writing \(t=re^{i\phi}\), it is
\begin{equation}
 x\longmapsto[r,(e^{i\phi},\theta_\gamma(x))]
 \in\operatorname{Cone}(\partial D_0\times S^1),
 \label{cusp-cone-q1-cusp-formula}
\end{equation}
and every point of \(W\) maps to \(c_0\).
\end{enumerate}
The fourth smash power is a global based map
\begin{equation}
 q_1^{\wedge4}:X^{\wedge4}\longrightarrow\mathcal C^{\wedge4}.
 \label{cusp-cone-S24-map}
\end{equation}
Using \(X\cong_{\mathrm{diff}}S^6\), its source is
\begin{equation}
 X^{\wedge4}\cong(S^6)^{\wedge4}\cong S^{24}.
 \label{cusp-cone-S24-source}
\end{equation}
On four synchronized radial fibres, the canonical product-to-smash quotient
maps restrict to embeddings on the source fibre product and the interior
target-product cross-section.  Under those embeddings,
\cref{cusp-cone-S24-map} restricts exactly to
\begin{equation}
 \chi_{\gamma,T}^{\times4}:F_T^4\longrightarrow C_T^4,\qquad
 C_T=\mathbb R/\ell_T\mathbb Z.
 \label{cusp-cone-S24-fibre-map}
\end{equation}
With its Euclidean \(\ell^2\)-product flat quotient metric,
\begin{equation}
 \operatorname{sys}(C_T^4)=\ell_T
 =T+H-B+O(T^{-1}),\qquad
 (C_T^4,0)\longrightarrow(\mathbb R^4,0)
 \label{cusp-cone-S24-R4-limit}
\end{equation}
in pointed local-isometric, pointed smooth, and pointed
Gromov--Hausdorff topology.

The target homology is
\begin{equation}
 \widetilde H_k(\mathcal C;\mathbb Z)\cong
 \begin{cases}
  \mathbb Z,&k=2,3,\\
  0,&\text{otherwise},
 \end{cases}
 \label{cusp-cone-C-homology}
\end{equation}
and
\begin{equation}
 \widetilde H_k(\mathcal C^{\wedge4};\mathbb Z)\cong
 \begin{cases}
  \mathbb Z,&k=8,12,\\
  \mathbb Z^4,&k=9,11,\\
  \mathbb Z^6,&k=10,\\
  0,&\text{otherwise}.
 \end{cases}
 \label{cusp-cone-C-smash-homology}
\end{equation}
\end{theorem}

\begin{proof}
The compact spaces \(P\times S^1,C_1,C_2\) are glued along closed boundary
tori by homeomorphisms, producing a compact orientable three-manifold
\(\mathcal M\) with boundary \(\partial D_0\times S^1\).  Attaching the cone
along this closed boundary gives a compact Hausdorff finite CW space
\(\mathcal C\).

The overlap identity
\[
 \chi_\gamma\circ h_j=h_j^{S^1}\circ\overline\chi_j
\]
proves that the first two pieces of \(q_1\) agree.  On \(\partial D_0\), the
cone base is glued to the remaining boundary of \(P\times S^1\), so the
first and third pieces agree.  If \(U\) is a neighborhood of \(c_0\), its
inverse image in the cone contains \([0,\delta)\times T^2\) for some
\(\delta>0\), by compactness of the cone base.  Therefore \(q_1^{-1}(U)\)
contains \(f^{-1}(|t|<\delta)\), a neighborhood of \(W\).  This proves
continuity at the new stratum; continuity elsewhere follows from the input
maps and the closed pasting lemma.

Every regular circle quotient and both finite Seifert targets are
surjective, and \(W\) maps onto the cone vertex.  Thus \(q_1\) is a
continuous surjection from compact \(X\) to Hausdorff \(\mathcal C\), hence
a quotient map.  It is based, so its fourth smash power exists.  Iterating
\(S^m\wedge S^n\cong S^{m+n}\) gives \cref{cusp-cone-S24-source}.

The base point lies in \(W\), while each radial regular fibre is disjoint
from \(W\).  Hence \(F_T^4\) and its product-circle image avoid the closed
source and target fat wedges.  The canonical quotient maps from the source
product and the indicated interior target-product cross-section to their
smash products restrict to embeddings.  Under those embeddings, the induced
smash map restricts to \(\chi_{\gamma,T}^{\times4}\), which itself retains
its displayed nontrivial torus kernel.
The quotient-metric calculation in
\cref{cusp-spec-global-circle-circumference} gives \(\ell_T\).
The product lattice is \(\ell_T\mathbb Z^4\), so the radius-\(R\) argument,
including the shortest-representative proof that the embedded Euclidean ball
is the whole quotient ball, in the proof of \cref{cusp-cone-S12-route} proves
\cref{cusp-cone-S24-R4-limit}.

For homology, let \(\rho_0,\rho_1,\rho_2\) be the boundary loops of the pair
of pants, with \(\rho_0+\rho_1+\rho_2=0\), and let \(\eta\) be the regular
target-circle class.  With the boundary-loop and circle orientations used
in the finite-filling presentation, the meridian relations are
\begin{equation}
 3\rho_1-\eta=0,\qquad4\rho_2+\eta=0.
 \label{cusp-cone-Seifert-meridians}
\end{equation}
Reversing the circle orientation reverses both displayed \(\eta\)-signs
and changes none of the following groups.  Thus
\begin{equation}
 H_1(\mathcal M;\mathbb Z)
 =\frac{\mathbb Z\langle\eta,\rho_1,\rho_2\rangle}
 {\langle3\rho_1-\eta,\;4\rho_2+\eta\rangle}
 \cong\mathbb Z\langle g\rangle,
 \label{cusp-cone-M-H1}
\end{equation}
where
\begin{equation}
 \rho_0=-g,\qquad\rho_1=4g,\qquad
 \rho_2=-3g,\qquad\eta=12g.
 \label{cusp-cone-M-boundary-map}
\end{equation}
In particular,
\(H_1(\partial\mathcal M)=\mathbb Z\langle\rho_0,\eta\rangle
\to H_1(\mathcal M)\) is surjective.  The orientable three-manifold
\(\mathcal M\) has one torus boundary, \(H_3(\mathcal M)=0\), and
\(\chi(\mathcal M)=\frac12\chi(\partial\mathcal M)=0\); hence
\(H_2(\mathcal M;\mathbb Q)=0\).  To obtain the integral statement without
suppressing torsion, Poincare--Lefschetz duality and the universal
coefficient theorem give
\[
 H_2(\mathcal M;\mathbb Z)
 \cong H^1(\mathcal M,\partial\mathcal M;\mathbb Z)
 \cong \operatorname{Hom}
 \bigl(H_1(\mathcal M,\partial\mathcal M;\mathbb Z),\mathbb Z\bigr),
\]
because \(H_0(\mathcal M,\partial\mathcal M;\mathbb Z)=0\).  Thus
\(H_2(\mathcal M;\mathbb Z)\) is free; its rank is zero, so it vanishes.

The cone attachment is the homotopy cofiber of
\(\partial\mathcal M\hookrightarrow\mathcal M\).  Since the boundary
inclusion is a cofibration,
\[
 \widetilde H_k(\mathcal C;\mathbb Z)
 \cong H_k(\mathcal M,\partial\mathcal M;\mathbb Z).
\]
Poincare--Lefschetz duality and
\(H^0(\mathcal M)=H^1(\mathcal M)=\mathbb Z\),
\(H^2(\mathcal M)=H^3(\mathcal M)=0\), give
\cref{cusp-cone-C-homology}.  The homology is free, so the smash Kunneth
theorem says that four factors of degree \(2\) or \(3\) produce ranks
\(\binom4j\) in degree \(8+j\), proving
\cref{cusp-cone-C-smash-homology}.
\end{proof}

\begin{definition}[The three radial target-slice embeddings]
\label{cusp-cone-target-slice-embeddings}
Put \(r_T=e^{-2\pi T}\) and let
\[
 \vartheta_T:\mathbb R/\mathbb Z\longrightarrow
 C_T=\mathbb R/\ell_T\mathbb Z,\qquad [a]\longmapsto[\ell_Ta].
\]
For \(y_1,y_2\in Q_T\), define
\[
 j_T(y_1,y_2)=
 \bigl[[r_T,1,\iota_T^{-1}(y_1)],
       [r_T,1,\iota_T^{-1}(y_2)]\bigr]_{\wedge}
 \in\mathcal A^{\wedge2}.
\]
For \(c_1,\ldots,c_4\in C_T\), define
\[
 k_T^{\mathcal B}(c_1,\ldots,c_4)=
 \bigl[[r_T,1,\vartheta_T^{-1}(c_1)],\ldots,
       [r_T,1,\vartheta_T^{-1}(c_4)]\bigr]_{\wedge}
 \in\mathcal B^{\wedge4},
\]
and, using the cusp-cone piece of \(\mathcal C\),
\[
 k_T^{\mathcal C}(c_1,\ldots,c_4)=
 \bigwedge_{j=1}^4[r_T,(1,\vartheta_T^{-1}(c_j))]
 \in\mathcal C^{\wedge4}.
\]
These are embeddings because every displayed factor lies away from the smash
basepoint, so equality in the smash quotient is equality in every factor;
compact-to-Hausdorff injectivity then gives a topological embedding.
\end{definition}

\begin{proposition}[Intrinsic expansion versus ambient collapse]
\label{cusp-cone-metric-scope}
Let \(D_{12}\), \(D_{24}^{\mathcal B}\), and
\(D_{24}^{\mathcal C}\) be arbitrary fixed metrics compatible with the
topologies of \(\mathcal A^{\wedge2}\), \(\mathcal B^{\wedge4}\), and
\(\mathcal C^{\wedge4}\), respectively.  In their corresponding ambient
metrics, the radial target
slices satisfy
\begin{equation}
 \operatorname{diam}_{\mathrm{amb}}
 \bigl(j_T(Q_T^2)\bigr)\longrightarrow0,\qquad
 \operatorname{diam}_{\mathrm{amb}}
 \bigl(k_T^{\mathcal B}(C_T^4)\bigr)\longrightarrow0,\qquad
 \operatorname{diam}_{\mathrm{amb}}
 \bigl(k_T^{\mathcal C}(C_T^4)\bigr)\longrightarrow0.
 \label{cusp-cone-ambient-collapse}
\end{equation}
By contrast, in the intrinsic quotient metrics,
\begin{equation}
 \operatorname{diam}(Q_T^2)\geq\tfrac12(T-C)\longrightarrow\infty,\qquad
 \operatorname{diam}(C_T^4)\geq\tfrac12\ell_T\longrightarrow\infty.
 \label{cusp-cone-intrinsic-expansion}
\end{equation}
Thus \cref{cusp-cone-S12-R4-limit,cusp-cone-S24-R4-limit} assert
intrinsic, flat pointed limits of varying slices, not Sormani--Wenger
intrinsic-flat convergence and not an ambient pointed \(\mathbb R^4\) limit
inside one fixed compact metric target.
\end{proposition}

\begin{proof}
Every neighborhood of the relevant cone vertex contains every sufficiently
small radial cross-section, uniformly in all compact angular and torus
coordinates.  In a compatible metric, every such cross-section eventually
lies in an arbitrarily small ball about the vertex.  Products followed by
the smash quotient preserve this uniform convergence, proving
\cref{cusp-cone-ambient-collapse}.

For a flat torus with systole \(s\), choose a shortest primitive period
\(\lambda\).  The class of \(\lambda/2\) has distance \(s/2\) from zero:
for every period \(\mu\), the nonzero period \(\lambda-2\mu\) has length at
least \(s\).  Hence the diameter is at least \(s/2\).
Equations \cref{cusp-cone-native-systole,cusp-cone-S24-R4-limit} prove
\cref{cusp-cone-intrinsic-expansion}.
\end{proof}

\begin{proposition}[Exact almost-complex obstruction on \(S^{12}\)]
\label{cusp-cone-S12-almost-complex-obstruction}
Let \(M\) be a smooth manifold homeomorphic to \(S^{12}\).  There is no
bundle endomorphism \(J:TM\to TM\) satisfying \(J^2=-\operatorname{Id}_{TM}\).
This conclusion has two independent proofs which use only the homotopy type
of \(M\): the original Borel--Serre reduced-power calculation at the single
prime \(5\), and Bott integrality.  For the standard embedded sphere there is
also an independent stable-tangent Pontryagin-class contradiction.

These statements concern an unstabilized complex structure on the actual
rank-twelve tangent bundle.  A stable complex structure on
\(TS^{12}\oplus\varepsilon^r_{\mathbb R}\), a complex structure confined to
an open subset or to fibres of a map, or an endomorphism which ceases to be a
bundle endomorphism along a singular set is not an almost-complex structure
on \(TS^{12}\) and therefore lies outside the hypothesis.
\end{proposition}

\begin{proof}
Assume that \(J\) exists and put \(E=(TM,J)\), a complex vector bundle of
rank six.  Orient \(M\) by the interleaved complex frame
\[
 (v_1,Jv_1,\ldots,v_6,Jv_6),
\]
write \([M]_J\) for the resulting fundamental class, and choose
\(u_J\in H^{12}(M;\mathbb Z)\) by
\(\langle u_J,[M]_J\rangle=1\).  Since \(M\) has the integral cohomology of
\(S^{12}\),
\begin{equation}
 c_1(E)=c_2(E)=c_3(E)=c_4(E)=c_5(E)=0.
 \label{cusp-cone-S12-lower-Chern-zero}
\end{equation}
The top-Chern/Euler identity uses exactly the orientation just fixed, so
Poincare--Hopf gives
\begin{equation}
 c_6(E)=e(E_{\mathbb R})=e(TM)=2u_J.
 \label{cusp-cone-S12-top-Chern-two}
\end{equation}
Thus reversing a previously chosen orientation can change the displayed
generator but cannot change the coefficient two.

We first specialize the original calculation of Borel and Serre
\cite[\S12, Proposition~12.4, Lemma~12.5, and Proposition~12.6,
p.~430; \S15, Proposition~15.1, p.~434]{BorelSerre1953}.
They write \(C_{2j}\) for the universal class now denoted \(c_j\).
For \(j=6\), the admissible prime is exactly
\[
 p=5<6,\qquad 5\nmid6,
\]
and their identity in \(H^*(BU(6);\mathbb F_5)\) reads
\begin{equation}
 \overline c_6
 =6^{-1}\mathcal P_5^1(\overline c_2)
   +Q(\overline c_1,\ldots,\overline c_5).
 \label{cusp-cone-S12-Borel-Serre-specialization}
\end{equation}
Here \(6^{-1}=1\) in \(\mathbb F_5\).  To retain the exact coefficient rather
than suppressing the lower-class polynomial, pass to the splitting space,
write \(x_1,\ldots,x_6\) for the degree-two Chern roots, put
\(s_r=\sum_i x_i^r\), and use
\(\mathcal P_5^1(x_i)=x_i^5\) and the Cartan formula.  The resulting universal
identity is
\begin{align}
 \mathcal P_5^1(\overline c_2)
 ={}&\sum_{i<j}(x_i^5x_j+x_ix_j^5)=s_5s_1-s_6 \notag\\
 ={}&6\overline c_6+\overline c_1^4\overline c_2
 -\overline c_1^3\overline c_3
 -4\overline c_1^2\overline c_2^2
 +\overline c_1^2\overline c_4
 +7\overline c_1\overline c_2\overline c_3
 -\overline c_1\overline c_5 \notag\\
 &\quad+2\overline c_2^3-6\overline c_2\overline c_4
 -3\overline c_3^2.
 \label{cusp-cone-S12-P1-exact}
\end{align}
Consequently, for the lower-Chern ideal
\(I=(\overline c_1,\ldots,\overline c_5)\),
\begin{equation}
 \mathcal P_5^1(\overline c_2)
 \equiv6\overline c_6=\overline c_6\pmod I,
 \label{cusp-cone-S12-P1-lower-ideal}
\end{equation}
which verifies the exact top-class coefficient in
\cref{cusp-cone-S12-Borel-Serre-specialization}.  Pulling back along the
classifying map of \(E\), using \cref{cusp-cone-S12-lower-Chern-zero} and
naturality of \(\mathcal P_5^1\), gives
\[
 \overline{c_6(E)}=\mathcal P_5^1(\overline{c_2(E)})
 =\mathcal P_5^1(0)=0.
\]
But \cref{cusp-cone-S12-top-Chern-two} gives
\(\overline{c_6(E)}=2\overline u_J\ne0\) in
\(H^{12}(M;\mathbb F_5)\), the first contradiction.

For the second contradiction, choose an orientation-preserving homotopy
equivalence \(f:S^{12}\to M\).  Bott integrality says that the sixth Chern
class of every complex bundle on \(S^{12}\) is divisible by
\((6-1)!=120\)
\cite[Chapter~20, \S9, Theorem~9.6 and Corollary~9.8,
p.~308]{Husemoller1994}; see also
\cite[\S6]{KonstantisParton2018}.  The coefficient can be checked without
normalization ambiguity: Newton's sixth identity and
\cref{cusp-cone-S12-lower-Chern-zero} give
\[
 \operatorname{ch}_6(f^*E)
 =\frac{s_6(f^*E)}{6!}
 =-\frac{c_6(f^*E)}{5!}.
\]
The integral Bott generator therefore forces
\(c_6(f^*E)=120k f^*u_J\) for some \(k\in\mathbb Z\), whereas
\cref{cusp-cone-S12-top-Chern-two} gives
\(c_6(f^*E)=2f^*u_J\).  Thus \(120k=2\), which has no integral solution.

Finally take \(M=S^{12}\) with its standard smooth embedding in
\(\mathbb R^{13}\).  There is an explicit stable trivialization
\begin{equation}
 TS^{12}\oplus\varepsilon^1_{\mathbb R}
 \xrightarrow{\ \cong\ }\varepsilon^{13}_{\mathbb R},
 \qquad (x;v,a)\longmapsto(x;v+ax),
 \label{cusp-cone-S12-stable-trivialization}
\end{equation}
so \(p_3(TS^{12})=0\).  On the other hand,
\[
 E_{\mathbb R}\otimes_{\mathbb R}\mathbb C
 \cong E\oplus\overline E,
 \qquad
 p_3(E_{\mathbb R})=-c_6(E\oplus\overline E).
\]
Equation \cref{cusp-cone-S12-lower-Chern-zero} and
\(c_6(\overline E)=c_6(E)\) give the exact relation
\begin{equation}
 p_3(TS^{12})=-2c_6(E)=-4u_J,
 \label{cusp-cone-S12-Pontryagin-contradiction}
\end{equation}
contradicting stable triviality in the torsion-free group
\(H^{12}(S^{12};\mathbb Z)\).  This is the Pontryagin--Euler proof recorded
independently in
\cite[Lemma~4.1 and Theorem~4.2]{KonstantisParton2018}.

Every contradiction above requires a globally defined \(J\) on the actual
bundle \(TM\).  Stabilizing that bundle, deleting points, restricting to
regular fibres, or allowing \(J\) to fail at a singular locus removes that
hypothesis rather than producing an exception to it.
\end{proof}

\begin{theorem}[Sharp obstructions to stronger interpretations]
\label{cusp-cone-obstructions}
The sphere routes above have the following exact limitations.
\begin{enumerate}
\item Let \(r\geq1\), and let \(N\subset\Lambda^{\oplus r}\) be a
primitive kernel defining a connected fibrewise torus quotient over the
full regular base, or over any finite cover of that base.  Equip the quotient
with the norm induced by the unrescaled marked Euclidean metric on the
original fibres.  If its quotient systole tends to infinity on the radial
cusp, then
\begin{equation}
 \operatorname{rank}(\Lambda^{\oplus r}/N)\leq r.
 \label{cusp-cone-rank-bound}
\end{equation}
Thus two copies cannot give a full-base four-torus quotient.  The
\(S^{12}\) map in \cref{cusp-cone-S12-route} evades the obstruction only by
localizing to the cusp and collapsing the exterior.

\item No nontrivial finite base cover, including one extending with
branching over the original singular values, makes \(K_\Lambda\) a
sub-local system over the full regular base.  The local cyclic change
\(t=s^M\) used for the auxiliary degree-\(M^2\) cover cannot extend as a
connected cover of \(B\cong\mathbb P^1\) branched only at \(p_0\).

\item There is no map from \(X^2\), \(X\wedge X\), \(X^4\), or
\(X^{\wedge4}\) to a fixed \(T^4\) whose restriction retains all four
target character classes of the displayed product quotient.  There is no
map to a fixed \(S^4\) whose restriction is a degree-one collapse of
\(Q_T^2\), or of the four-circle target, for even one regular fibre.  The
conical targets evade these obstructions because the fibre character and
orientation classes do not extend as global target classes.

\item Cartesian products do not produce the named spheres:
\begin{equation}
 X^2\simeq S^6\times S^6\not\simeq S^{12},\qquad
 X^4\simeq(S^6)^4\not\simeq S^{24}.
 \label{cusp-cone-products-not-spheres}
\end{equation}
The two- and four-fold synchronized fibre products have real dimensions
\(10\) and \(18\), respectively, and generally have singular completed
fibres.  Joins give
\begin{equation}
 X*X\cong S^{13},\qquad X^{*4}\cong S^{27}.
 \label{cusp-cone-join-dimensions}
\end{equation}

\item Since \(X\cong S^6\) is simply connected, every connected unbranched
finite cover of \(X\), \(X^r\), \(X\wedge X\cong S^{12}\), or
\(X^{\wedge4}\cong S^{24}\) is trivial.  The local degree-\(M^2\)
unspooling cover is not the restriction of a nontrivial connected
unbranched cover of any of these compact sources.  Branched or normalized
compactifications are different objects and are not constructed here.

\item By \cref{cusp-cone-S12-almost-complex-obstruction}, no smooth manifold
homeomorphic to \(S^{12}\) admits an almost-complex structure.  Nor does
\(S^{24}\), with any differentiable structure, by the Borel--Serre theorem
\cite[\S15, Proposition~15.1 and the following remark,
p.~434]{BorelSerre1953}.  Thus the smash sources cannot be promoted to
compact complex spheres.
\end{enumerate}
\end{theorem}

\begin{proof}
Every element of \(K_\Lambda^{\oplus r}\) has uniformly bounded physical
length on the radial cusp.  If one survived in the quotient lattice, it
would be a nonzero period of bounded quotient norm.  Hence
\begin{equation}
 K_\Lambda^{\oplus r}\subset N.
 \label{cusp-cone-kernel-contained}
\end{equation}
For the full monodromy group,
\(A_1\widehat w=\widehat u-\widehat w\) and invariance force every copy of
\(\widehat u\) into \(N\).  Therefore
\[
 L_\Lambda^{\oplus r}
 =\mathbb Z\langle\widehat u,\widehat w,\widehat\delta\rangle^{\oplus r}
 \subset N,
\]
which proves \cref{cusp-cone-rank-bound}.

The same conclusion holds after a finite base cover.  Let \(H\) be its
finite-index subgroup and put \(H_*=A_2A_1=I+N_*\), \(N_*^2=0\).
Every finite-index subgroup contains \(H_*^n\) for some \(n>0\), while
\begin{equation}
 H_*^n\widehat w
 =\widehat w+n(\widehat u-\widehat\delta).
 \label{cusp-cone-finite-index-orbit}
\end{equation}
Equations
\cref{cusp-cone-kernel-contained,cusp-cone-finite-index-orbit} give
\(n\widehat u\in N\).  Primitivity gives \(\widehat u\in N\), proving the
same rank bound.  The orbit also gives
\(H_*^nK_\Lambda\ne K_\Lambda\) for every \(n>0\), so no finite cover of
the regular base preserves the native kernel.  A compact cover branched
over the singular values restricts to such a cover.  Finally, a connected
finite cover of \(\mathbb P^1\) branched only at \(p_0\) would restrict to
a connected unbranched finite cover of
\(\mathbb P^1\setminus\{p_0\}\cong\mathbb C\), hence has degree one.

For a fixed \(T^4\) target, \(H^1(X^r;\mathbb Z)=0\), as does the first
cohomology of the smash spheres.  A global map therefore pulls every target
character class back to zero.  On a regular fibre the quotient homomorphisms
inject the target character lattice into the character lattice of the
source torus, giving four independent nonzero pullbacks, a contradiction.

For the fixed \(S^4\) obstruction, let
\(\omega_T\in H^2(Q_T;\mathbb Z)\) be the orientation class.  Since
\(\pi_T^*:H^1(Q_T)\to H^1(F_T)\) injects two independent classes and torus
cohomology is their exterior algebra,
\begin{equation}
 (\pi_T\times\pi_T)^*(\omega_T\times\omega_T)
 \ne0\quad\text{in }H^4(F_T^2;\mathbb Z).
 \label{cusp-cone-fibre-H4-nonzero}
\end{equation}
A degree-one collapse \(Q_T^2\to S^4\) pulls the orientation generator
back to \(\pm(\omega_T\times\omega_T)\).  But
\[
 H^4(X^2;\mathbb Z)=H^4(S^{12};\mathbb Z)=0,
\]
so no global extension has this restriction.  The four-circle case is the
same: its orientation class pulls back to the exterior product of four
independent circle characters, while
\(H^4(X^4)=H^4(S^{24})=0\).

The Kunneth theorem gives \(H_6(X^2)=\mathbb Z^2\), whereas
\(H_6(S^{12})=0\), and gives nonzero groups of ranks \(4,6,4\) in degrees
\(6,12,18\) for \(X^4\), whereas \(S^{24}\) has none.  The real dimension
of an \(r\)-fold synchronized fibre product is
\(6r-2(r-1)=4r+2\), giving \(10\) and \(18\).  The join formula
\(S^m*S^n\cong S^{m+n+1}\) proves
\cref{cusp-cone-join-dimensions}.  The remaining assertions follow from
simple connectivity, \cref{cusp-cone-S12-almost-complex-obstruction}, and
the cited Borel--Serre theorem for \(S^{24}\).
\end{proof}

\paragraph{Canonicity and exact scope.}
The lattice \(K_\Lambda=\ker(M_0-I)=\operatorname{im}(M_0-I)\) and the
punctured-cusp quotient are intrinsic to the marked cusp monodromy.
The compact map \(q_2\) additionally depends on the chosen cusp coordinate,
closed cusp disc, and the decision to collapse its outer boundary and the
exterior.  The identifications
\(X^{\wedge2}\cong S^{12}\) and \(X^{\wedge4}\cong S^{24}\) depend on a
based choice of \(X\cong S^6\).  None of these choices is claimed to be
holomorphic or canonical in the algebraic category.

The map \(q_2^{\wedge2}\) is the direct two-copy construction from the
native \(T^2\times\mathbb R^2\) cusp: it is globally defined on the smash
sphere, and its radial regular-fibre quotients are \(Q_T\times Q_T\), which
converge pointed local-isometrically, smoothly, and Gromov--Hausdorff to
\(\mathbb R^4\).  It is not a torus quotient over the full
regular base.  The map \(q_{1,\mathrm{cusp}}^{\wedge4}\) is the parallel
cusp-supported construction with target
\(\operatorname{Susp}(T^2)^{\wedge4}\).  The stronger map
\(q_1^{\wedge4}\) uses the quotient over every regular and finite-filled
fibre before the cusp cone is attached; four copies are minimal for that
full-base all-expanding group quotient.
Neither construction gives a compact complex \(S^{12}\) or \(S^{24}\), an
ambient Euclidean limit in a fixed compact metric target, or a claim about
the interacting Yang--Mills spectrum.

The argument does not rule out a separately constructed branched or
toroidal compactification of the auxiliary finite covers after additional
branching away from \(p_0\).  Such a construction would have to specify its
global base cover, normalize and resolve every pulled-back singular fibre,
track the lifted metric, and determine the topology of the resulting total
space.  No such smooth compact sphere construction follows from the present
data.
